\documentclass[11pt]{article}
\usepackage[a4paper,margin=25mm]{geometry}
\usepackage{hyperref}
\usepackage{booktabs}
\title{From Contribution Growth to Maintenance Capacity:\\A Longitudinal Study Protocol for Headscale}
\author{Research protocol derived from the Headscale coursework}
\date{13 October 2026}
\begin{document}
\maketitle
\paragraph{Status.} This is a living study protocol updated after the motivating
pilot, not a preregistration or submission-ready manuscript. Completed work includes
release-pinned Git analysis, core-alias and timestamp sensitivity, all-state issue
collection, the complete PR metadata census, a 76-PR process pilot with timelines,
evidence extraction for 100 issues, complete submitted-review and lifecycle
collection for 2,166 PRs, and exploratory first-episode merge-hazard models with
role/experience, exposure-window and elapsed-time diagnostics. Component-matched
role sensitivity uses complete file lists for 1,318 merged PRs. Independent semantic
outcome coding and explanatory mechanism validation remain incomplete. No acceptance
claim or causal conclusion is made.

\paragraph{Target contribution.} Establish whether growth in contribution breadth
translates into distributed maintenance responsibility, and whether concentrated
maintenance is associated with different response and resolution outcomes. Test
both bottleneck and efficient-specialisation explanations. The upstream-funded
alternative implementation is a distinctive case context, not proof that funding
causes either outcome. Initial target: an empirical study suitable for consideration
by Journal of Systems and Software, Information and Software Technology or
Empirical Software Engineering; venue suitability must be reassessed after results.

\paragraph{Research questions.}
\begin{enumerate}
\item RQ1: How do contribution, review and merge concentration evolve in fixed
calendar windows, and do they move together?
\item RQ2: What relationship can be established between prior maintenance
concentration, workload and subsequent response or disposition outcomes?
\item RQ3: Which observable closure and responsibility mechanisms change the
interpretation of repository-level maintenance indicators?
\end{enumerate}

\paragraph{Novelty gate.} Contribution ownership and response latency are existing
research topics. Reusing HHI or analysing a popular repository is not a novelty
claim. Position against ownership and review-aware ownership studies, integrator
surveys, CODEOWNERS research, time-to-first-response work and recent PR-workflow
mortality research. Advance only if validated process data
and triangulated mechanisms provide a distinct contribution; otherwise expand the
comparative design or narrow the manuscript to a measurement-validity study.

\paragraph{Evidence and provenance.} Preserve the coursework and its closed-bug
cache. Pin each Git endpoint by tag and commit SHA; record the official release
publication timestamp separately from commit timestamps. Use an independent bare
repository so working-tree edits cannot enter measurements. Retain API pages,
retrieval timestamps, integrity hashes and completion status. Treat historical
issue labels and state as unknown where only collection-time observations exist.
PR reviews, comments, merge data and issue timelines are distinct data sources.
Record coverage and missingness for each before analysis.

\paragraph{Measurements.} Primary time windows are complete calendar quarters,
with monthly or 90-day windows as sensitivity analyses. Partial windows are labelled
and excluded from direct volume comparisons. Authorship is contribution activity,
not knowledge. The minimum number of contributors accounting for 50\% of activity
is named literally, not called a validated bus factor. Separate production Go,
tests, generated code, documentation and other artifacts. Track confirmed file
renames and manually audit splits and moves. Compare name-based and validated
identity mappings. Measure human response and review directly from timestamped
events, excluding author self-comments and bots under a documented rule. Define
merge concentration from merger identities, not commit authors.

\paragraph{Defects and survival outcomes.} Begin with all issues, including open
ones. Determine defect eligibility using labels plus human validation. Distinguish
verified fix, duplicate, unsupported configuration, cannot reproduce, deliberate
non-fix, and unknown closure. Closure and repair are separate outcomes. Link repair
PRs/commits where supported by explicit evidence. Include open issues as right
censored at the observation cutoff, handle reopened episodes explicitly, and treat
non-fix closure as a competing event. Report time-to-closure descriptively even
when fix linkage is unavailable, without calling it MTTR. Historical claims must
not use future events or collection-time states as if known at release time.

\paragraph{Validation design.} Use a fixed-seed stratified random sample across
periods, closure types, automated/human actors and latency tails. Two independent
coders should validate defect status and closure outcomes, reconcile disagreements
and report agreement with uncertainty. A single-coder pilot must be labelled as
such. Record exclusions and sampling probabilities. Prepare interview questions
and public-evidence coding, but contact participants only after explicit user
authorization and the applicable research-ethics requirements are established.
Public FOSDEM testimony is contextual evidence, not an independent confirmation
of repository-derived conclusions.
On 14 October the author deferred human coding. The two-reviewer and explicitly
declared single-reviewer options remain available; neither has been carried out.
Other manuscript preparation continues without assigning human labels or
reporting repair prevalence. Before a mechanism claim, trace a component-boundary
observation to an explicit maintenance decision and subsequent treatment, retaining
unclear and contradictory cases. This is a post-pilot explanatory test, not a
preregistered hypothesis.

\paragraph{Inference plan.} First establish coverage, identities and event
definitions. Report effect sizes and uncertainty; plot distributions and survival
curves rather than only means. Model lagged concentration with workload, backlog,
change size, contributor experience, component, issue type and available severity
controls. Account for calendar trends, dependence and clustered observations.
Use competing-risk or cause-specific models when justified; assess proportional
hazards before using a Cox interpretation. Avoid inferring causality from temporal
co-movement. Choose models and comparison-project sample sizes after checking
feasibility and the estimand, not after selecting favourable results.

\paragraph{External validation gate.} A single case can support a carefully bounded
mechanism study with triangulation. Claims about open-source projects generally
require a transparent comparative sample. Select projects using documented ecosystem,
language, age, activity and governance criteria, without selecting on the outcome.
Exclude incomparable label/review histories or report them explicitly. If no
credible comparison can be assembled, keep conclusions specific to Headscale.

\paragraph{Exploratory identity sensitivity, 14 October 2026.}
After inspecting missing review accounts, add a fixed-population sensitivity
for fractional reviewer concentration. Retain merged PRs with an attributed
eligible review. Apply the same cutoff, open intervals and final-disposition
boundaries to unattributed submitted records. Derive conservative bounds from
known-credit floors and the maximum additional identity capacity; do not describe
relaxed endpoints as sharp bounds or confidence intervals. Report all affected
quarters and conditional account scenarios, including changed role orderings.
This addition follows observation of the data and is not preregistered.
Deleted or informal review and known-account alias uncertainty remain uncovered.

\paragraph{Exploratory component validation, 15 October 2026.}
The execution plan in \texttt{repair\_validation\_plan.json} was saved before
running the checks. From the fixed issue sample, select all final-closure commit
pointers with an exact merged-PR match and a locally available official commit.
This source-availability rule selects issues 2121 and 2181 from five pointers.
Compare each repair commit with its first parent, holding the Go toolchain,
dependency manifests and regression oracle constant within the pair. Retain
normal-case controls. One oracle follows the reported duplicate-name rename
setup; the other uses the exact upstream ACL regression test file. Require
executed target tests to fail before and pass after the repair, with controls
passing in both variants. Compilation failures, absent tests and timeouts are
inconclusive. Report component behavior separately from original-deployment
reproduction, independent human outcome coding and population repair prevalence.
This is purposive post-pilot triangulation, not a preregistered experiment.

\paragraph{Exploratory interpretation audits, 15 October 2026.}
Joint account participation is now compared on the full merged-PR population,
separating self-merging from attributed other-account submitted review. Collected
approval states and exact matches to the collected head revision are secondary
record-linkage diagnostics, not historical policy-enforcement claims. A separate
calendar audit reproduces the frozen linear models, adds a four-degree-of-freedom
natural cubic spline and restricts elapsed-duration risk-set comparisons within
creation-quarter baseline strata. All executed spline previews are disclosed;
the plans are post-pilot. Stratum support, excluded intervals and late events
remain explicit. Different conditional comparisons do not establish a causal
effect or a formal difference between model coefficients.

\paragraph{Preliminary update.} Legacy Git analysis at v0.28.0 gives 2,803
non-merge, non-bot commits after the initial tag, 238 case-insensitive author names,
48.16\% top-author share and author HHI 0.2602. At v0.29.5 these become 3,401,
259, 54.78\% and 0.3202. The newly segmented cross-author edit interval median
is 13.95 days, versus 21.27 in the preceding long segment. These measurements use
unequal windows and are descriptive; the interval is not measured coordination
waiting. They motivate testing divergent explanations rather than retaining the
coursework's monotonic ``scaling crisis'' narrative. New issue outcomes have not
yet been established by this Git update.

\paragraph{Submission readiness gate.} Required: validated data definitions;
complete or explicitly bounded event coverage; reproducible pinned analysis;
appropriate uncertainty and robustness; related-work positioning; triangulated
interpretation or comparative validation; a complete manuscript with bounded
claims; checked figures, tables and references. A Linux port, larger dataset or
manuscript skeleton alone does not meet this gate.
\end{document}
